\section{Cutoff and the discrete update clock}
\label{rcut:sec:cutoff}

Section~\ref{rcut:sec:dynamics} placed the median mixing time on a
logarithmic scale.  We now show that, once the chain has a fixed positive
amount of common mass with equilibrium, any additional positive
logarithmic lag makes its total-variation distance vanish.  Applying this
implication on either side of the median proves continuous-time cutoff;
a comparison of update clocks then gives the discrete-time result.

Fix $\beta\in[0,1/2)$.  Throughout this section, the disorder belongs to
the simultaneous good events $\mathcal G_n$ from
Section~\ref{rcut:sec:dynamics}.  In particular, $\|J\|\le K$, the spectral
gap is at least $\gamma>0$, and the constants in
Equation~\eqref{rcut:eq:gradient-propagation} are uniform on these events.
We use its notation $\kappa>0$, $C_{\mathrm{dyn}}$, and $r_n$.
For each fixed $D_*>0$, Equation~\eqref{rcut:eq:gradient-propagation} applies whenever
$0\le s\le d\le D_*\log n$, and its deterministic remainder satisfies
\begin{equation}
 n^p r_n\longrightarrow0\qquad\text{for every fixed }p>0.
 \label{rcut:eq:cutoff-remainder}
\end{equation}
We retain this uniformity throughout: unless stated otherwise, every
limit is uniform over $\mathcal G_n$.  Constants may depend on the fixed
inverse temperature and on the fixed parameters in the relevant
statement.

\subsection{Common mass and local fields}

If two laws share a fixed positive amount of mass, their means can
be separated only to the extent permitted by their variances.  The
following elementary estimate makes this precise.  We first apply it to
local fields; later we apply it to bounded test functions after
semigroup smoothing.

\begin{lemma}[Means under measures with common mass]
\label{rcut:lem:common-mass}
Let $\nu$ and $\pi$ be probability measures on a finite set, and suppose
$\|\nu-\pi\|_{\mathrm{TV}}\le1-\varepsilon$, where
$\varepsilon>0$.  For every real function $H$ on that set,
\begin{equation}
 |\nu H-\pi H|
 \le \frac{\sqrt{\Var_\nu H}+\sqrt{\Var_\pi H}}
              {\sqrt{\varepsilon}}.
 \label{rcut:eq:common-mass-means}
\end{equation}
\end{lemma}

\begin{proof}
Take the common part of the measures,
$\rho(x)=\min\{\nu(x),\pi(x)\}$, whose mass is
$q=1-\|\nu-\pi\|_{\mathrm{TV}}\ge\varepsilon$.
We compare each mean to the mean under the normalized common part.
Since $\rho\le\nu$, Cauchy--Schwarz gives
\[
 \left|\nu H-\frac{\rho H}{q}\right|
 =\frac1q\left|\sum_x\rho(x)(H(x)-\nu H)\right|
 \le\sqrt{\frac{\Var_\nu H}{q}}.
\]
The same calculation for $\pi$ bounds the other difference from the
common mean.  The triangle inequality, followed by the lower bound on
the common mass, proves the assertion.
\end{proof}

Write $h_i(x)=(Jx)_i$ for the local field.  Its coefficient vector has
Euclidean norm at most $K$.  Spin inversion gives $\mu h_i=0$ in the
zero-field model.  Thus the concentration estimate for linear
observables applies to every local field.  Common mass will also
control its transient mean, yielding the following tail bound at each
deterministic observation time.

\begin{lemma}[Fields at deterministic observation times]
\label{rcut:lem:cutoff-field-tails}
Fix $\varepsilon\in(0,1)$, $\delta>0$, and $D_*>0$.
Suppose $d_c(t)\le1-\varepsilon$.  There is a deterministic sequence
$\tau_n=\tau_n(\varepsilon,\delta,D_*)$ such that
\begin{equation}
 \sup_x\P_x\left[\max_{1\le i\le n}|h_i(X_z)|>\delta\log n\right]
 \le\tau_n,
 \qquad t\le z\le D_*\log n,
 \label{rcut:eq:cutoff-field-tails}
\end{equation}
and $n^p\tau_n\to0$ for every fixed $p>0$.
\end{lemma}

\begin{proof}
First control the means.  Contraction of total variation implies that
the law $\nu_{x,z}$ of $X_z$ from any point $x$ has common mass at least
$\varepsilon$ with $\mu$ for every $z\ge t$.
Proposition~\ref{rcut:prop:linear-concentration} bounds
$\Var_{\nu_{x,z}}h_i$ uniformly in $i,x,z$, and the equilibrium
variance is at most $K^2/\gamma$.  Lemma~\ref{rcut:lem:common-mass} therefore
gives a constant $B_\varepsilon$ such that
\begin{equation}
 |\E_x h_i(X_z)|\le B_\varepsilon
 \qquad(t\le z\le D_*\log n).
 \label{rcut:eq:cutoff-field-means}
\end{equation}

Next control the deviations from these means.  For every fixed
$\theta>0$, Proposition~\ref{rcut:prop:linear-concentration} bounds both
centered exponential moments by a common constant $M_{\theta,K}$.
Apply the exponential bound to both tails of each field and then take
a union bound over the coordinates.  Together with
Equation~\eqref{rcut:eq:cutoff-field-means}, this gives
\begin{equation}
 \P_x\left[\max_i|h_i(X_z)|>\delta\log n\right]
 \le 2M_{\theta,K}e^{\theta B_\varepsilon}
             n^{1-\theta\delta}.
 \label{rcut:eq:cutoff-field-chernoff}
\end{equation}
For a prescribed $p>0$, first fix $\theta$ with
$\theta\delta>p+1$, and then let $n$ tend to infinity.  Define
$\tau_n$ as the supremum of the displayed probabilities over the
indicated parameters and good disorders.  The same sequence therefore
has every asserted inverse-power decay rate.  The uniformity is over
deterministic times.  In the smoothing argument we will integrate the
resulting expectation bounds, so a simultaneous field event along the
whole sample path is unnecessary.
\end{proof}

\subsection{A positive logarithmic lag finishes mixing}

The passage from a nontrivial total-variation overlap to mixing after a
positive logarithmic lag is related to the concentration-based cutoff
mechanism of Salez~\cite[Theorem~5 and Lemma~9]{Salez2024}.  Here we
prove the required implication through variances of bounded smoothed
tests and their common mass with equilibrium.

We now prove the implication needed for cutoff.  The common-mass
hypothesis will compare expectations of a smoothed test once its
variance is small under both the transient law and equilibrium.
The main task is the transient variance: gradient propagation will
bound the martingale energy before the observation time by the energy
in a later interval of positive logarithmic length.

\begin{proposition}[Positive-lag smoothing]
\label{rcut:prop:positive-lag}
Fix $\varepsilon\in(0,1)$, $\alpha>0$, and $D_*>0$, and put
$s=\alpha\log n$.  Uniformly over all times $t\ge0$ satisfying
\[
 d_c(t)\le1-\varepsilon,
 \qquad t+s\le D_*\log n,
\]
one has $d_c(t+s)\to0$.
\end{proposition}

\begin{proof}
Fix a starting point $x$ and a real test function with
$\|f\|_\infty\le1$.  We first express the transient variance as an
integrated martingale energy.  Set $T=t+s$ and
\[
 F_u=S_{T-u}f,\qquad
 G(u)=\E_x\Gamma(F_u)(X_u),\qquad 0\le u\le T.
\]
Here $F_u(X_u)$ is the backward martingale with terminal value
$f(X_T)$.  Differentiating its second moment and using
$(\partial_u+L)F_u=0$ gives the exact identity
\begin{equation}
 \Var_x(F_r(X_r))=2\int_0^rG(u)\,du,
 \qquad 0\le r\le T.
 \label{rcut:eq:cutoff-variance-budget}
\end{equation}
Since $\|F_u\|_\infty\le1$, this implies
\begin{equation}
 \int_0^T G(u)\,du\le\frac12,
 \qquad
 \Var_x(S_sf(X_t))=2\int_0^tG(u)\,du.
 \label{rcut:eq:cutoff-budget-consequences}
\end{equation}

The total energy is thus bounded.  To show that
little of it lies before the observation time, fix deterministic
$u\le t$ and $z\in[t+s/2,t+s]$.  Apply
Equation~\eqref{rcut:eq:gradient-propagation} starting at $X_u$, with
duration $T-u$ and comparison time $z-u$.  These parameters lie within
the fixed horizon, and the estimate holds for every possible value of
$X_u$ with the same constants.  Integrating against its law and using
the Markov property yields
\[
 \E_x U(F_u)(X_u)
 \le C_{\mathrm{dyn}}e^{-\kappa(z-u)}
                  \E_x U(F_z)(X_z)+n r_n.
\]
The bound on the remainder uses $U(F_z)\le n$: each half-difference of a
function bounded by one has magnitude at most one.  Combining this
with $\Gamma\le2U$ and the separation $z-u\ge s/2$ gives
\begin{equation}
 G(u)\le 2C_{\mathrm{dyn}}n^{-\kappa\alpha/2}
                  \E_x U(F_z)(X_z)+2n r_n.
 \label{rcut:eq:cutoff-early-late}
\end{equation}

We still need to replace the late unweighted gradient by the late
martingale energy.  The field-tail estimate supplies the required
lower bound on the heat-bath weights.  Choose a fixed $\delta$ with
$0<\delta<\kappa\alpha/4$, and let $\tau_n$ be the sequence in
Lemma~\ref{rcut:lem:cutoff-field-tails}.  At every configuration,
\[
 w_i=1-x_i\tanh h_i
 \ge 1-|\tanh h_i|
 =\frac{2}{1+e^{2|h_i|}}\ge e^{-2|h_i|}.
\]
On the event $\max_i|h_i(X_z)|\le\delta\log n$ this gives
$U(F_z)(X_z)\le n^{2\delta}\Gamma(F_z)(X_z)$.
On the complementary event use $U(F_z)\le n$.  These are pointwise
bounds, so they remain valid even though the fields and $F_z$ depend
on the same configuration.  Taking expectations gives
\begin{equation}
 \E_xU(F_z)(X_z)\le n^{2\delta}G(z)+n\tau_n.
 \label{rcut:eq:cutoff-field-gradient}
\end{equation}
Define the positive exponent and the deterministic remainder
\[
 \xi=\frac{\kappa\alpha}{2}-2\delta>0,
 \qquad R_n=2n r_n+2C_{\mathrm{dyn}}n\tau_n.
\]
Equations~\eqref{rcut:eq:cutoff-early-late} and
\eqref{rcut:eq:cutoff-field-gradient} imply
\begin{equation}
 G(u)\le 2C_{\mathrm{dyn}}n^{-\xi}G(z)+R_n
 \quad\text{for }u\le t,\quad t+s/2\le z\le t+s.
 \label{rcut:eq:cutoff-energy-comparison}
\end{equation}
The choice of the field threshold leaves a strictly positive decay
exponent, and $R_n$ decreases faster than every inverse power of $n$.
We have therefore compared each early energy to every energy in the
late interval, with uniform error.  Average
Equation~\eqref{rcut:eq:cutoff-energy-comparison} over the interval of
$z$ of length $s/2$, and then integrate over $0\le u\le t$.
The total energy budget in
Equation~\eqref{rcut:eq:cutoff-budget-consequences} gives
\begin{equation}
 \Var_x(S_sf(X_t))
 \le 4C_{\mathrm{dyn}}\frac{t}{s}n^{-\xi}+2tR_n
 \le 4C_{\mathrm{dyn}}\frac{D_*}{\alpha}n^{-\xi}
            +2D_*\log n\,R_n=:V_n.
 \label{rcut:eq:cutoff-late-variance}
\end{equation}
Thus $V_n\to0$ uniformly in $f,x,t$ and in the good disorder, proving
the required transient variance bound.  The integrations concern
deterministic expectations of finite-state semigroup expressions;
their measurability and nonnegativity justify interchanging the two
integrals by Tonelli's Theorem.  No random observation time or union
over a continuum of field events enters this step.

It remains to control the equilibrium variance and compare the two
means.  The uniform spectral gap from
Proposition~\ref{rcut:prop:uniform-gap} gives $\frac{d}{dr}\Var_\mu(S_rf)\le-2\gamma\Var_\mu(S_rf)$.
Integrating yields
\begin{equation}
 \Var_\mu(S_sf)\le e^{-2\gamma s}\Var_\mu(f)
                        \le n^{-2\gamma\alpha}.
 \label{rcut:eq:cutoff-equilibrium-smoothing}
\end{equation}
Both variance bounds now apply to the same observable.  Use
Lemma~\ref{rcut:lem:common-mass} for the law of $X_t$ from $x$ and
for $\mu$, with observable $S_sf$.  The common-mass hypothesis is
exactly the assumed distance bound at the observation time, and the
two expectations are $S_{t+s}f(x)$ and $\mu f$.  Hence
\[
 |S_{t+s}f(x)-\mu f|
 \le \varepsilon^{-1/2}
          \bigl(\sqrt{V_n}+n^{-\gamma\alpha}\bigr).
\]
Taking one half of the supremum over $x$ and $\|f\|_\infty\le1$
proves Proposition~\ref{rcut:prop:positive-lag}.
\end{proof}

Positive-lag smoothing turns the logarithmic median scale into cutoff.
At the median there is enough common mass to apply the proposition.
If such common mass appeared a fixed fraction earlier, the same
proposition would force mixing before the median, which is impossible.

\begin{theorem}[Continuous-time cutoff]
\label{rcut:thm:continuous-cutoff}
Let $t_n^{\rm med}=\inf\{t\ge0:d_c(t)\le1/2\}$ be the continuous-time median
mixing time.  For every fixed $b\in(0,1)$,
\begin{equation}
 d_c((1-b)t_n^{\rm med})\longrightarrow1,
 \qquad d_c((1+b)t_n^{\rm med})\longrightarrow0,
 \label{rcut:eq:continuous-cutoff-sides}
\end{equation}
uniformly over $\mathcal G_n$.
\end{theorem}

\begin{proof}
Proposition~\ref{rcut:prop:mixing-scale} gives fixed constants $c,C>0$
such that $c\log n\le t_n^{\rm med}\le C\log n$ for all sufficiently large $n$,
and $d_c(t_n^{\rm med})=1/2$.  Fix $b\in(0,1)$ and set $\alpha=bc/2$.
Fix the horizon constant $D_*>(1+b)C$ so that both applications of
positive-lag smoothing lie within its range.  First apply
Proposition~\ref{rcut:prop:positive-lag} at $t=t_n^{\rm med}$.  This gives
$d_c(t_n^{\rm med}+\alpha\log n)\to0$.  Since
$\alpha\log n\le bt_n^{\rm med}/2$, monotonicity yields the upper conclusion.

For the lower side, argue by contradiction while preserving uniformity
in the disorder.  If the conclusion failed uniformly, there would be a subsequence
of dimensions, a good disorder at each of those dimensions, and a
fixed $\varepsilon_0>0$ such that
$d_c((1-b)t_n^{\rm med})\le1-\varepsilon_0$.  Proposition~\ref{rcut:prop:positive-lag}
would then give vanishing distance at time
\[
 (1-b)t_n^{\rm med}+\alpha\log n
 \le t_n^{\rm med}-\frac{bc}{2}\log n<t_n^{\rm med}.
\]
This contradicts the definition of $t_n^{\rm med}$, since the distance at every
time strictly below $t_n^{\rm med}$ exceeds $1/2$.
\end{proof}

\subsection{A clock comparison with fixed holding}

We next transfer cutoff to the clock that counts individual update
attempts.  Poisson concentration gives the lower side directly.  For
the upper side we first extract a fixed holding probability, then apply
the uniform Poisson-to-binomial comparison proved in
Lemma~\ref{lem:input-clock}.  This is the same finite-chain comparison
used in the main argument; below we verify all its hypotheses for the
present median scale.

\begin{proposition}[Cutoff for single-site update attempts]
\label{rcut:prop:discrete-cutoff}
For every fixed $b\in(0,1)$, uniformly over $\mathcal G_n$,
\begin{equation}
 d_d\bigl(\lfloor(1-b)nt_n^{\rm med}\rfloor\bigr)\longrightarrow1,
 \qquad
 d_d\bigl(\lceil(1+b)nt_n^{\rm med}\rceil\bigr)\longrightarrow0.
 \label{rcut:eq:discrete-cutoff-sides}
\end{equation}
Here each discrete step is one uniformly selected heat-bath update
attempt, including attempts that leave the configuration unchanged.
\end{proposition}

\begin{proof}
The exact clock relation is $S_t=e^{nt(P-I)}$, so continuous time
averages discrete update attempts against a Poisson law.  For the
lower bound, put
$k=\lfloor(1-b)nt_n^{\rm med}\rfloor$ and $T=(1-b/2)t_n^{\rm med}$, and let
$N\sim\operatorname{Poisson}(nT)$.  Convexity of total variation and
monotonicity of the discrete distance give
\begin{equation}
 d_c(T)\le\sum_{j\ge0}\P(N=j)d_d(j)
                \le\P(N<k)+d_d(k).
 \label{rcut:eq:clock-lower-transfer}
\end{equation}
The Poisson mean exceeds $k$ by at least $(b/2)nt_n^{\rm med}$.
Chebyshev's inequality therefore implies $\P(N<k)=O((nt_n^{\rm med})^{-1})=o(1)$ uniformly.
Theorem~\ref{rcut:thm:continuous-cutoff}, with parameter $b/2$, gives
$d_c(T)\to1$; Equation~\eqref{rcut:eq:clock-lower-transfer} proves the
discrete lower bound.

For the upper bound, we must check the fixed holding condition used
by the clock-comparison lemma.  The operator-norm bound implies
\[
 \frac1n\sum_i h_i(x)^2=\frac{\|Jx\|^2}{n}\le K^2.
\]
At least half the fields have magnitude at most $\sqrt2K$.  When spin
$i$ is selected, its chance of being retained is
$e^{x_ih_i}/(2\cosh h_i)\ge(1+e^{2|h_i|})^{-1}$.
Averaging this lower bound over the selected coordinate gives
\begin{equation}
 P(x,x)\ge h_0:=\frac{1}{2(1+e^{2\sqrt2K})}>0.
 \label{rcut:eq:clock-uniform-holding}
\end{equation}
Choose $q=1-h_0/2\in(0,1)$ once for all $n$ and all good
disorders, and write
\[
 P=(1-q)I+q R,
 \qquad R=\frac{P-(1-q)I}{q}.
\]
Subtracting this fixed holding component leaves a Markov kernel
$R$ with stationary measure $\mu$.  In terms of its powers, the two
clocks become
\begin{equation}
 S_t=e^{q nt(R-I)},\qquad
 P^M=\sum_{j=0}^M\binom Mjq^j(1-q)^{M-j}R^j.
 \label{rcut:eq:clock-thinning}
\end{equation}

We now verify convergence of every fixed shifted Poisson average,
as required by Lemma~\ref{lem:input-clock}.  Set
$M=\lceil(1+b)nt_n^{\rm med}\rceil$.  For any $x\in\Omega_n$ and any
real function $f$ with $\|f\|_\infty\le1$, define
$u_j=(R^jf)(x)-\mu f$, so that $|u_j|\le2$.
For fixed $y\in\R$, Equation~\eqref{rcut:eq:clock-thinning} identifies
the Poisson average with mean $m=q M+y\sqrt M$ as
\[
 \E\bigl[u_{\operatorname{Poisson}(m)}\bigr]
 =S_{m/(q n)}f(x)-\mu f.
\]
Uniformly over the good disorders, $t_n^{\rm med}$ is between fixed positive
multiples of $\log n$.  Hence
\[
 \frac{m}{q n}
 =\frac Mn+\frac{y\sqrt M}{q n}
 \ge(1+b/2)t_n^{\rm med}
\]
for all sufficiently large $n$: the correction is
$O_y(\sqrt{\log n/n})$, whereas $(b/2)t_n^{\rm med}$ is of order $\log n$.
Theorem~\ref{rcut:thm:continuous-cutoff} and monotonicity bound this Poisson
average in absolute value by $2d_c((1+b/2)t_n^{\rm med})=o(1)$, uniformly over
$x$, $f$, and the good disorder.  Also $M$ tends uniformly to infinity.
All hypotheses of Lemma~\ref{lem:input-clock} are now uniform over
the indicated family, so its binomial average tends uniformly to zero.
Equation~\eqref{rcut:eq:clock-thinning} identifies that average with
$P^Mf(x)-\mu f$.
Taking one half of the supremum over $x\in\Omega_n$ and real
$f$ with $\|f\|_\infty\le1$ proves the upper bound in
Equation~\eqref{rcut:eq:discrete-cutoff-sides}.
\end{proof}

\begin{proof}[Completion of the proof of Theorem~\ref{rcut:thm:main}]
It remains to express discrete cutoff as the asserted ratio of mixing
times and then return to disorder probability.  Fix
$\varepsilon\in(0,1/2)$ and $\eta>0$.  Choose once
\[
 0<b<\frac{\eta}{2+\eta},
 \qquad\text{so that}\qquad
 \frac{1+b}{1-b}<1+\eta.
\]
Proposition~\ref{rcut:prop:discrete-cutoff} implies, for all sufficiently
large $n$ and every disorder in $\mathcal G_n$,
\[
 t_{n,\beta,J}(1-\varepsilon)>
                \lfloor(1-b)nt_n^{\rm med}\rfloor,
 \qquad
 t_{n,\beta,J}(\varepsilon)\le
                \lceil(1+b)nt_n^{\rm med}\rceil.
\]
The first bound ensures a positive denominator and controls the
integer rounding in the ratio.  Hence
\[
 \frac{t_{n,\beta,J}(\varepsilon)}
      {t_{n,\beta,J}(1-\varepsilon)}
 \le\frac{\lceil(1+b)nt_n^{\rm med}\rceil}
          {\lfloor(1-b)nt_n^{\rm med}\rfloor+1}
 \le\frac{1+b}{1-b}+\frac{1}{(1-b)nt_n^{\rm med}}<1+\eta
\]
for all sufficiently large $n$, uniformly on $\mathcal G_n$.
The last step uses $t_n^{\rm med}\ge c\log n$.
Since $\P_J(\mathcal G_n)\to1$, the probability of the event in
Theorem~\ref{rcut:thm:main} is at most $\P_J(\mathcal G_n^c)$ for large $n$,
and therefore tends to zero.  This last step uses only the probability
of the good event: all preceding estimates hold for each fixed disorder,
uniformly over the starting state and the test function.  No coupling
of disorders across dimensions is needed.
\end{proof}
